% from main.tex
%%%%%%%%%%  scan idx 298  |  folio 285  %%%%%%%%%%
désigne le rang) :
\[
\begin{tikzcd}[column sep=tiny]
H^0(Y,\mathbb Z)\arrow[r,"i"]\arrow[d]
&K'(Y)_{\naif}\arrow[r]\arrow[d]
&K'(Y)\arrow[r,"\Rg"]\arrow[d]
&H^0(Y,\mathbb Z)\arrow[d]\\
H^0(X,\mathbb Z)\arrow[r,"i"]
&K'(X)_{\naif}\arrow[r]
&K'(X)\arrow[r,"\Rg"]
&H^0(X,\mathbb Z).
\end{tikzcd}
\]
\textsuperscript{*}On verra dans la suite du Séminaire (V 3.9.1 et VI) que \(K'(X)_{\naif}\) peut être muni d'une structure plus riche que celle d'algèbre augmentée vers \(H^0(X,\mathbb Z)\), à savoir d'une structure de \(\lambda\)-algèbre augmentée vers \(H^0(X,\mathbb Z)\) considéré comme anneau binômial (V 2.8), cette structure étant compatible avec les homomorphismes image inverse. Deligne [5] a prouvé récemment que \(K'(X)\) peut être muni d'une structure analogue (du moins si \(X\) est de caractéristique \(p\geq0\)) et que l'homomorphisme canonique
\[
 K'(X)_{\naif}\longrightarrow K'(X)
\]
est un homomorphisme de \(\lambda\)-algèbres augmentées, répondant ainsi par l'affirmative à une conjecture de Grothendieck (XIV 1).\textsuperscript{*}

\textbf{2.9. Comparaison de \(K'\) et \(K'_{\naif}\).} Soit \((X,A)\) un topos annelé. Il n'y a pas de raison en général pour que l'homomorphisme canonique
\[
 K'(X)_{\naif}\longrightarrow K'(X)
\]
soit injectif ou surjectif. Cependant, les «~critères d'existence de résolutions globales~» de Exp. II montrent que c'est un isomorphisme dans les cas suivants :
\begin{enumerate}[label=(\roman*)]
\item \(X\) est le topos zariskien d'un schéma divisoriel, séparé ou noethérien (et \(A=\mathcal O_X\)) (II 2.2.9);
\item \(X\) est le topos défini par un espace compact et \(A\) est un faisceau mou (II 2.3.2);
\item \(X\) est un topos de Stein dont l'objet final est quasi-compact et \(A=\mathcal O_X\) (II 2.4.5).
\end{enumerate}

